% GENERATED FILE. Edit canonical lessons, then run npm run generate:books.
% canonical-source-hash: fnv1a64:d889591d8de838f5
% canonical-lessons: RU-C229-sekretar, RU-C229-bukhgalter, RU-C229-menedzher, RU-C229-programmist, RU-C229-prepodavatel

\chapter{Secretary, Accountant, Manager, Programmer, Lecturer}
\label{ch:ru-secretary-accountant-manager-programmer-lecturer}
\hlchaptermodality{229}

\begin{chapteropening}
\textbf{By the end of this chapter:} \emph{I can say a secretary, an accountant, a manager, a programmer and a lecturer.}

Complete the last lesson of chapter 229: I can say a secretary, an accountant, a manager, a programmer and a lecturer.
\end{chapteropening}

\section[\texorpdfstring{\ru{секретарь} (sekretár)}{секретарь (sekretár)}]{\ru{секретарь} (sekretár) — a secretary}

\label{lesson:RU-C229-sekretar}



\begin{quote}
Before the new one: say the Russian for a meaning, then the Russian for a first name.
\end{quote}

\subsection*{You'll want to know: \ru{секретарь}}
\textbf{\ru{секретарь}} (\emph{sekretár}) — \textquotedblleft{}a secretary\textquotedblright{}.

\begin{cousinweb}
\textbf{\ru{секретарь}} — known from 1614; it came through German \emph{Sekretär} or Polish \emph{sekretarz} from Latin \emph{sēcrētārius}, \textquotedblleft{}a keeper of secrets\textquotedblright{} (Vasmer).
\end{cousinweb}

\begin{grammarlens}[title={What you've built}]
One more word for this chapter.
\end{grammarlens}

\subsection*{Your turn}
\begin{itemize}
\raggedright
  \item \emph{Say it:} \emph{sekretár}
  \item \emph{Say it:} \emph{sekretár}, once more
  \item \emph{Say it:} say \emph{sekretár}
  \item \emph{Recall:} say the Russian for a meaning, then the Russian for a first name, then say \emph{sekretár} again
\end{itemize}

\begin{tcolorbox}[breakable,colback=teal!4,colframe=teal!35!black,title={Before you move on}]
What does \ru{секретарь} mean? (\textquotedblleft{}A secretary\textquotedblright{}.) Say it once more.
\end{tcolorbox}
\section[\texorpdfstring{\ru{бухгалтер} (bukhgálter)}{бухгалтер (bukhgálter)}]{\ru{бухгалтер} (bukhgálter) — an accountant}

\label{lesson:RU-C229-bukhgalter}



\begin{quote}
Before the new one: say the Russian for a first name, then the Russian for a secretary.
\end{quote}

\subsection*{You'll want to know: \ru{бухгалтер}}
\textbf{\ru{бухгалтер}} (\emph{bukhgálter}) — \textquotedblleft{}an accountant\textquotedblright{}.

\begin{cousinweb}
\textbf{\ru{бухгалтер}} — from Peter the Great's time; it comes from German \emph{Buchhalter}, \textquotedblleft{}book-keeper\textquotedblright{} (Vasmer).
\end{cousinweb}

\begin{grammarlens}[title={What you've built}]
One more word for this chapter.
\end{grammarlens}

\subsection*{Your turn}
\begin{itemize}
\raggedright
  \item \emph{Say it:} \emph{bukhgálter}
  \item \emph{Say it:} \emph{bukhgálter}, once more
  \item \emph{Say it:} say \emph{bukhgálter}
  \item \emph{Recall:} say the Russian for a first name, then the Russian for a secretary, then say \emph{bukhgálter} again
\end{itemize}

\begin{tcolorbox}[breakable,colback=teal!4,colframe=teal!35!black,title={Before you move on}]
What does \ru{бухгалтер} mean? (\textquotedblleft{}An accountant\textquotedblright{}.) Say it once more.
\end{tcolorbox}
\section[\texorpdfstring{\ru{менеджер} (ménedzher)}{менеджер (ménedzher)}]{\ru{менеджер} (ménedzher) — a manager}

\label{lesson:RU-C229-menedzher}



\begin{quote}
Before the new one: say the Russian for a secretary, then the Russian for an accountant.
\end{quote}

\subsection*{You'll want to know: \ru{менеджер}}
\textbf{\ru{менеджер}} (\emph{ménedzher}) — \textquotedblleft{}a manager\textquotedblright{}.

\begin{cousinweb}
\textbf{\ru{менеджер}} comes from English \emph{manager}, from \emph{manage}, Italian \emph{maneggiare}, \textquotedblleft{}to handle\textquotedblright{}, from Latin \emph{manus}, \textquotedblleft{}hand\textquotedblright{}.
\end{cousinweb}

\begin{grammarlens}[title={What you've built}]
One more word for this chapter.
\end{grammarlens}

\subsection*{Your turn}
\begin{itemize}
\raggedright
  \item \emph{Say it:} \emph{ménedzher}
  \item \emph{Say it:} \emph{ménedzher}, once more
  \item \emph{Say it:} say \emph{ménedzher}
  \item \emph{Recall:} say the Russian for a secretary, then the Russian for an accountant, then say \emph{ménedzher} again
\end{itemize}

\begin{tcolorbox}[breakable,colback=teal!4,colframe=teal!35!black,title={Before you move on}]
What does \ru{менеджер} mean? (\textquotedblleft{}A manager\textquotedblright{}.) Say it once more.
\end{tcolorbox}
\section[\texorpdfstring{\ru{программист} (programmíst)}{программист (programmíst)}]{\ru{программист} (programmíst) — a programmer}

\label{lesson:RU-C229-programmist}



\begin{quote}
Before the new one: say the Russian for an accountant, then the Russian for a manager.
\end{quote}

\subsection*{You'll want to know: \ru{программист}}
\textbf{\ru{программист}} (\emph{programmíst}) — \textquotedblleft{}a programmer\textquotedblright{}.

\begin{cousinweb}
\textbf{\ru{программист}} is \emph{\ru{программа}}, \textquotedblleft{}a program\textquotedblright{}, with \emph{-\ru{ист}}, \textquotedblleft{}one who works at\textquotedblright{}. \emph{\ru{программа}} goes back to Greek \emph{prógramma}, \textquotedblleft{}a public notice\textquotedblright{}, from \emph{prográphein}, \textquotedblleft{}to write before\textquotedblright{}.
\end{cousinweb}

\begin{grammarlens}[title={What you've built}]
One more word for this chapter.
\end{grammarlens}

\subsection*{Your turn}
\begin{itemize}
\raggedright
  \item \emph{Say it:} \emph{programmíst}
  \item \emph{Say it:} \emph{programmíst}, once more
  \item \emph{Say it:} say \emph{programmíst}
  \item \emph{Recall:} say the Russian for an accountant, then the Russian for a manager, then say \emph{programmíst} again
\end{itemize}

\begin{tcolorbox}[breakable,colback=teal!4,colframe=teal!35!black,title={Before you move on}]
What does \ru{программист} mean? (\textquotedblleft{}A programmer\textquotedblright{}.) Say it once more.
\end{tcolorbox}
\section[\texorpdfstring{\ru{преподаватель} (prepodavátel)}{преподаватель (prepodavátel)}]{\ru{преподаватель} (prepodavátel) — a lecturer, a teacher}

\label{lesson:RU-C229-prepodavatel}



\begin{quote}
Before the new one: say the Russian for a manager, then the Russian for a programmer.
\end{quote}

\subsection*{You'll want to know: \ru{преподаватель}}
\textbf{\ru{преподаватель}} (\emph{prepodavátel}) — \textquotedblleft{}a lecturer, a teacher\textquotedblright{}.

\begin{cousinweb}
\textbf{\ru{преподаватель}} is \emph{\ru{преподавать}}, \textquotedblleft{}to teach\textquotedblright{}, with \emph{-\ru{тель}}, \textquotedblleft{}one who does it\textquotedblright{}. The verb is \emph{\ru{пре}-} and \emph{\ru{по}-} with \emph{\ru{давать}}, \textquotedblleft{}to give\textquotedblright{}.
\end{cousinweb}

\begin{grammarlens}[title={What you've built}]
One more word for this chapter.
\end{grammarlens}

\subsection*{Your turn}
\begin{itemize}
\raggedright
  \item \emph{Say it:} \emph{prepodavátel}
  \item \emph{Say it:} \emph{prepodavátel}, once more
  \item \emph{Say it:} say \emph{prepodavátel}
  \item \emph{Recall:} say the Russian for a manager, then the Russian for a programmer, then say \emph{prepodavátel} again
\end{itemize}

\begin{tcolorbox}[breakable,colback=teal!4,colframe=teal!35!black,title={Before you move on}]
What does \ru{преподаватель} mean? (\textquotedblleft{}A lecturer, a teacher\textquotedblright{}.) Say it once more.
\end{tcolorbox}
